\chapter{SRM Marrakech-Safi : Préparer l'Oral d'Ingénieur d'État (Affectation Préfecture de Marrakech)}

\begin{tcolorbox}[colback=red!4!white,colframe=red!60!black,title=\textbf{Pourquoi ce chapitre (ajouté le 19 novembre 2026)}]
La SRM Marrakech-Safi recrute en 2026 des ingénieurs d'État en informatique (développement, systèmes d'information, génie logiciel), par le même organisme que la SRM Souss-Massa (Brain Management, opération « SRM-MS-2026 ») : écrit le 1\textsuperscript{er} novembre 2026 à Marrakech, puis entretien oral. Les postes sont affectés par zone géographique ; « Préfecture de Marrakech » désigne la zone d'affectation, l'employeur reste la SRM. Tout ce qui vaut pour la SRM Souss-Massa (loi 83-21, guichet unique, cinq piliers de l'entretien du chapitre 7, méthode STAR) vaut ici : seuls les faits de l'entreprise et de la région changent.
\end{tcolorbox}

\section{La SRM Marrakech-Safi en repères vérifiés}
\begin{itemize}
    \item \textbf{Démarrage :} contrat de gestion en vigueur le \textbf{1\textsuperscript{er} novembre 2024}, dans la première vague (EcoActu, 04/11/2024). Contrat de gestion de 30 ans.
    \item \textbf{Opérateurs remplacés :} la \textbf{RADEEMA} (Marrakech), la \textbf{RADEES} (Safi) et l'activité de distribution de l'\textbf{ONEE} dans la région.
    \item \textbf{Périmètre :} préfecture de Marrakech et provinces d'Al Haouz, Chichaoua, El Kelaâ des Sraghna, Essaouira, Rehamna, Safi et Youssoufia ; 251 communes ; environ \textbf{4,9 millions d'habitants} (4 892 393 au recensement de 2024, HCP).
    \item \textbf{Capital et actionnaires :} 100 MDH ; État 25 \%, groupement des collectivités territoriales 40 \%, ONEE 25 \%, Conseil régional 10 \% ; programme prévisionnel d'investissement de 28 MMDH annoncé en 2024 (EcoActu, 04/11/2024).
    \item \textbf{Directeur général :} M. Abdellah Ilhami, nommé par le conseil d'administration (annonce du 12/12/2024, MAP Express) ; aucun changement trouvé en 2026. Le président du conseil d'administration n'a pas été confirmé par une source : ne pas le citer.
    \item \textbf{Électricité (2025--2026) :} plus de 65 000 km de réseau, plus de 1,5 million de clients électricité, 34 postes sources, plus de 4 TWh distribués par an (H24info, 25/05/2026).
\end{itemize}

\section{Ce qui intéresse un ingénieur informatique}
\begin{itemize}
    \item \textbf{Systèmes déjà en place :} supervision \textbf{SCADA}, système d'information géographique (\textbf{SIG}) et gestion des pannes (\textbf{OMS}).
    \item \textbf{Systèmes prévus :} gestion avancée de la distribution (\textbf{ADMS}), comptage intelligent (\textbf{AMI}) et gestion des données de compteurs (\textbf{MDMS}) (H24info, 25/05/2026).
    \item \textbf{Relation client :} un plan d'action commercial 2026--2027 annoncé début novembre 2026 vise la modernisation des agences, la digitalisation des services et la réduction des délais de traitement (titre de presse, Aujourd'hui le Maroc, 05/11/2026 ; détail non lu).
    \item \textbf{Ce que cela veut dire à l'oral :} unifier les données clients et techniques héritées de trois opérateurs (RADEEMA, RADEES, ONEE), sécuriser des systèmes industriels critiques (segmentation du réseau SCADA, authentification forte, journalisation, loi 05-20), fiabiliser la facturation et ouvrir des services en ligne simples pour 4,9 millions d'habitants.
\end{itemize}

\section{L'eau : le grand chantier de la région}
\begin{itemize}
    \item \textbf{Transfert d'eau dessalée de Safi vers le Grand Marrakech :} station de dessalement de Safi réalisée par le groupe OCP ; environ 185 km de conduites, 3 stations de pompage, un réservoir de 10 000 m³, un débit de 3,2 m³ par seconde (près de 100 millions de m³ par an) ; environ 4,2 MMDH financés par l'État et le ministère de l'Intérieur ; maître d'ouvrage : la SRM Marrakech-Safi (Le Matin, 21/06/2026). La mise en service était annoncée pour fin juillet 2026 ; \textbf{aucune confirmation officielle de la mise en service n'a été trouvée au 19/11/2026} : en parler comme d'un projet en cours de mise en service.
    \item \textbf{Eau en milieu rural :} extension des réseaux aux douars des provinces de Safi et de Youssoufia ; plus de 760 000 habitants ruraux visés (Le360, 26/04/2025).
\end{itemize}

\section{La préfecture de Marrakech et la région}
\begin{itemize}
    \item \textbf{Wali de la région Marrakech-Safi et gouverneur de la préfecture de Marrakech :} M. Khatib El Hebil, nommé le 19 octobre 2025 (Le Matin, 19/10/2025), précédemment wali de l'Oriental ; il a succédé à M. Farid Chourak.
    \item \textbf{Rôle d'une préfecture :} selon l'article 145 de la Constitution, les walis et gouverneurs représentent le pouvoir central, veillent à l'application des lois, exercent le contrôle administratif des collectivités territoriales et coordonnent les services déconcentrés de l'État. Pour l'usager : CNIE et passeports (avec la DGSN), état civil, autorisations.
    \item \textbf{Préfecture de Marrakech :} 1 330 465 habitants (RGPH 2024), 15 communes ; chef-lieu de la région.
    \item \textbf{Séisme d'Al Haouz (8 septembre 2023) :} au 31 août 2026, 52 794 logements terminés, soit environ 91 \% du programme de relogement, et 5 126 chantiers encore en cours (Le360, 24/10/2026).
    \item \textbf{Mondial 2030 :} le Grand Stade de Marrakech est en transformation depuis février 2026 (stade dédié au football, toiture, capacité portée à environ 46 000 places), livraison annoncée fin 2028 (La Tribune de Marrakech, 23/08/2026) ; la LGV Kénitra--Marrakech est prévue pour fin 2029.
\end{itemize}

\section{Le numérique public à connaître}
\begin{itemize}
    \item \textbf{Loi 55-19} sur la simplification des procédures administratives et portail \textbf{Idarati} : nouvelle version le 18 septembre 2026 avec 874 services en ligne et 2 533 actes publiés (TelQuel, 21/09/2026).
    \item \textbf{État civil électronique} (loi 36-21) : registre numérique national, identifiant civil et social à la naissance, numérisation des registres papier demandée aux walis fin 2025 (Hespress FR, 27/10/2025).
    \item \textbf{Chikaya.ma} (réclamations des usagers) et \textbf{Watiqa} (extraits d'état civil en ligne).
    \item \textbf{Nouveau ministère} de l'Intelligence artificielle, de la Transformation numérique et de la Réforme de l'administration (M. Ahmed Rahhou, gouvernement du 16 octobre 2026) ; à l'Intérieur, M. Abdelouafi Laftit est reconduit.
\end{itemize}

\section{Questions probables et pistes de réponse}

{\small
\begin{longtable}{|>{\raggedright\arraybackslash\hspace{0pt}}p{5.2cm}|>{\raggedright\arraybackslash\hspace{0pt}}p{8.8cm}|}
\hline
\rowcolor{gray!20} \textbf{Question probable} & \textbf{Piste de réponse courte} \\ \hline
\endfirsthead
\hline
\rowcolor{gray!20} \textbf{Question probable} & \textbf{Piste de réponse courte} \\ \hline
\endhead

Que savez-vous de la SRM Marrakech-Safi ? & Trois faits datés : démarrage le 1\textsuperscript{er} novembre 2024 à la place de la RADEEMA, de la RADEES et de l'ONEE ; environ 4,9 millions d'habitants et 251 communes ; le transfert d'eau dessalée de Safi vers Marrakech. Puis le lien avec le poste : les systèmes SCADA, SIG et OMS, et les projets ADMS et comptage intelligent. \\ \hline

Pourquoi voulez-vous nous rejoindre ? & Un service essentiel (l'eau et l'électricité de toute une région), une entreprise jeune qui construit son système d'information à partir de trois anciens opérateurs, et la possibilité de voir l'effet direct de son travail pour les habitants. Ne jamais critiquer le privé. \\ \hline

Comment uniformiser les systèmes hérités de trois opérateurs ? & Cartographier l'existant, choisir un référentiel client unique, migrer par étapes avec des tests de rapprochement des données, garder un plan de retour en arrière, former les agents. \\ \hline

Comment protéger le SCADA ? & Séparer le réseau industriel du réseau bureautique, authentification forte, accès à distance contrôlé, journalisation et supervision, mises à jour planifiées, plan de continuité ; conformité à la loi 05-20 et aux directives de la DGSSI. \\ \hline

À quoi sert le comptage intelligent ? & Relever les compteurs à distance, détecter plus vite les fuites et les fraudes, facturer la consommation réelle, informer le client ; il faut en retour protéger les données personnelles (loi 09-08) et dimensionner la plateforme de données (MDMS). \\ \hline

Une panne touche un quartier de Marrakech un jour de forte chaleur : que faites-vous ? & Localiser avec l'OMS et le SIG, informer les équipes terrain et les clients (SMS, site), suivre le rétablissement, puis analyser la cause et documenter. \\ \hline

\end{longtable}
}

\begin{tcolorbox}[colback=yellow!5!white,colframe=orange!75!black,title=\textbf{À ne pas dire}]
\begin{itemize}
    \item Une date de mise en service du transfert Safi--Marrakech non confirmée.
    \item Le nom d'un président du conseil d'administration non vérifié.
    \item Des chiffres d'investissement 2026--2030 vus sur un seul site secondaire.
    \item Des sujets sportifs contestés (finale de la CAN 2025) : rester sur le Mondial 2030.
\end{itemize}
\end{tcolorbox}
